For a positive integer $n$, a {\it partition} of $n$ is a non-increasing sequence $\lambda = \lp \lambda_1, \dots, \lambda_\ell \rp$ of positive integers such that $|\lambda| := \lambda_1 + \dots + \lambda_\ell = n$. We write $\lambda \vdash n$ to denote that $\lambda$ is a partition of $n$, and we let $\ell = \ell\lp \lambda \rp$ be the number of parts of $\lambda$. Since the very beginning of partition theory, generating functions have been an essential tool. In the last century, the {\it circle method}
as developed first by Hardy and Ramanujan \cite{HR} has become an indispensable tool in studying the coefficients of these generating functions as well as many counting problems in additive combinatorics and number theory, including the Waring problem and weak Goldbach conjecture. Hardy and Ramanujan used the generating function for the {\it partition function} $p(n) := \#\{ \lambda \vdash n \}$, given for $|q|<1$ by
\begin{align*}
 \sum_{n \geq 0} p(n) q^n = \prod_{n \geq 1} \dfrac{1}{1 - q^n},
\end{align*}
in order to prove an asymptotic series for $p(n)$ as $n \to \infty$, whose main term yields
\begin{align*}
 p(n) \sim \dfrac{1}{4n\sqrt{3}} {\rm e}^{\pi\sqrt{\frac{2n}{3}}}.
\end{align*}

